% !TEX program = xelatex
% Шаблон научной статьи НИТУ «МИСиС»
% Заполните блоки ниже и скомпилируйте: xelatex article.tex

\documentclass{misisarticle}

% --- Классификационные коды (оставьте пустым {}, если не требуются) ---
\grnti{28.23.37, 28.23.25, 28.23.29}
\asjc{1702}
\oecd{1.02}

% --- Название ---
\title{Контекстно-зависимый классификатор услуг по текстам объявлений на основе трансформерной модели BERT}

% --- Авторы: \misisauthor{ФИО}{аффилиация}{город, страна}{email} ---
\misisauthor{А.~Р.~Тихомирова}{кафедра инженерной кибернетики НИТУ «МИСиС»}{Москва, Россия}{m2310563@edu.misis.ru}
\misisauthor{П.~П.~Научный}{кафедра инженерной кибернетики НИТУ «МИСиС»}{Москва, Россия}{supervisor@misis.ru}

% --- Аннотация и ключевые слова ---
\abstract{На онлайн-платформах для размещения объявлений продавец услуг указывает одну микрокатегорию, однако в описании может упоминать несколько — как самостоятельных, так и входящих в состав комплексной работы. Из-за этого объявления не попадают в релевантные микрокатегории, а пользователи не находят нужные услуги. Сложность состоит в том, что статус одной и той же услуги не фиксирован: он зависит от контекста, в котором она упомянута. В работе предлагается контекстно-зависимый классификатор на основе трансформерной модели BERT, дообученной на размеченном наборе объявлений с Avito. В отличие от существующих подходов, классификатор решает задачу многометочной классификации и одновременно предсказывает две характеристики: множество упоминаемых услуг и необходимость разделения объявления на несколько самостоятельных. Предложенный подход позволяет автоматически определять статус услуг и выделять самостоятельные из комплексных. Это даёт возможность платформам формировать дополнительные объявления в релевантных микрокатегориях, повышая видимость услуг для пользователей и снижая потери продавцов. 
}

\keywords{Классификация текстов, Multi-label классификация, Дообучение трансформера, Неструктурированный текст, Эмбеддинги, BERT
}

\begin{document}

\maketitle

% ============================================================
\section{Введение}

Первый абзац раздела начинается с красной строки (отступ 0{,}51~см). Текст выравнивается по ширине, шрифт Times New Roman, 10~pt.

Второй абзац. Ссылки на литературу оформляйте так~\cite{example2024}. Для нескольких источников:~\cite{example2024, book2023}.

% ============================================================
\section{Основная часть}

\subsection{Подраздел}

Подзаголовки второго уровня набираются \textit{курсивом}.

Пример маркированного списка:
\begin{itemize}
  \item первый пункт;
  \item второй пункт;
  \item третий пункт.
\end{itemize}

Пример формулы в строке: мера Жаккара $IoU = \dfrac{|A \cap B|}{|A \cup B|}$.

Нумерованная формула (ссылка в тексте: формула~\eqref{eq:f1}):
\begin{equation}
  F_1 = 2 \cdot \frac{P \cdot R}{P + R}
       = \frac{2TP}{2TP + FP + FN},
  \label{eq:f1}
\end{equation}
где $P$~--- precision, $R$~--- recall, $TP$, $FP$, $FN$~--- число истинно и ложно положительных/отрицательных детекций.

% --- Рисунок ---
\begin{figure}[htbp]
  \centering
  % Замените заглушку на свой файл:
  % \includegraphics[width=0.75\linewidth]{figures/figure}
  \fbox{\parbox[c][3cm][c]{0.9\columnwidth}{\centering Placeholder для рисунка}}
  \caption{Пример иллюстрации. Замените заглушку на свой файл из папки figures.}
  \label{fig:example}
\end{figure}

% --- Таблица ---
\begin{table}[htbp]
  \caption{Пример таблицы с результатами}
  \label{tab:example}
  \centering
  \begin{tabular}{lcc}
    \toprule
    \textbf{Метод} & \textbf{Precision} & \textbf{Recall} \\
    \toprule
    Метод~A & 0.85 & 0.78 \\
    Метод~B & 0.91 & 0.83 \\
    \bottomrule
  \end{tabular}
\end{table}

% ============================================================
\section{Заключение}

Кратко сформулируйте основные выводы работы.

% ============================================================
% Список литературы — вариант 1: вручную
\begin{thebibliography}{9}

\bibitem{example2024} Author A., Author B. Title of the paper. \textit{Journal Name}, vol.~1, no.~2, pp.~100--110, 2024.

\bibitem{book2023} Author C. \textit{Book Title}. Publisher, City, 2023.

\end{thebibliography}

\end{document}
